\documentclass[11pt]{article}

\usepackage[margin=1in]{geometry}
\usepackage{amsmath,amssymb,amsthm,mathtools}
\usepackage{microtype}
\usepackage{enumitem}
\usepackage{xcolor}
\usepackage{hyperref}
\usepackage{booktabs}
\usepackage{array}
\usepackage{longtable}
\usepackage{url}

\hypersetup{
  colorlinks=true,
  linkcolor=blue!55!black,
  citecolor=blue!55!black,
  urlcolor=blue!55!black,
  pdftitle={Polylogarithmic Genus Does Not Increase the Power of Constant-Width Polynomial-Size Circuits},
  pdfauthor={Grisha Pochuev}
}

\newtheorem{theorem}{Theorem}[section]
\newtheorem{lemma}[theorem]{Lemma}
\newtheorem{proposition}[theorem]{Proposition}
\newtheorem{corollary}[theorem]{Corollary}
\theoremstyle{definition}
\newtheorem{definition}[theorem]{Definition}
\newtheorem{remark}[theorem]{Remark}
\newtheorem{verification}[theorem]{Verification statement}

\newcommand{\ACC}{\mathrm{ACC}^{0}}
\newcommand{\ACm}[1]{\mathrm{AC}^{0}[#1]}
\newcommand{\genus}{\operatorname{genus}}
\newcommand{\bits}{\{0,1\}}
\newcommand{\cB}{\mathcal{B}}
\newcommand{\cK}{\mathcal{K}}
\newcommand{\cP}{\mathcal{P}}
\newcommand{\cR}{\mathcal{R}}
\newcommand{\ceil}[1]{\left\lceil #1\right\rceil}
\newcommand{\floor}[1]{\left\lfloor #1\right\rfloor}

\title{Polylogarithmic Genus Does Not Increase the Power of\\
Constant-Width Polynomial-Size Circuits\\
\large A Separator-Based Proof with Partial Lean Verification}
\author{Grisha Pochuev\\\small Independent researcher}
\date{11 August 2026 -- Version 6.0}

\setlength{\emergencystretch}{3em}
\begin{document}
\maketitle

\begin{abstract}
Allender, Datta, and Roy claimed that constant-width polynomial-size Boolean
circuits of polylogarithmic orientable genus compute exactly $\ACC$.  The
published proof of the upper bound was later found to contain a false
topological assertion, and Eric Allender subsequently listed the question as
open with a US\$1000 prize.  We give a different proof that does not use the
invalid handle-placement argument.

The central topological lemma is elementary once genus additivity is used in
the right way.  If a layered graph has $N$ vertices and orientable genus at
most $g$, then deleting at most
\[
  g\bigl(\ceil{\log_2 N}+1\bigr)
\]
whole layers makes the graph planar.  In every round we take one weighted
median layer from each currently nonplanar connected component.  There are at
most $g$ such components, and every nonplanar component surviving to the next
round has a canonical nonplanar parent of at least twice its size.  Hence no
nonplanar component survives more than $\ceil{\log_2 N}+1$ rounds.

For a constant-width circuit, these cuts yield only polylogarithmically many
macroblocks with a constant-size state interface.  Good macroblocks are
planar; bad macroblocks are single transitions.  We package all good block
output circuits into one polynomial-size planar family before invoking
Hansen's planar characterization, which forces one common modulus.  Finally,
a polylogarithmic sequence of relations on the fixed state space
$Q=\bits^w$ is composed in a constant number of rounds, each round explicitly
disjoining over only polynomially many intermediate state sequences.

A substantial part of the proof has also been machine checked in Lean 4.  In
the normalized fixed-width, fan-in-two source model used by the project, Lean
checks the canonical layer-planarization construction, macroblock semantics,
the simultaneous padding argument, finite-state relation composition, and the
construction of the final fixed-modulus target family.  The final checked
declaration is \nolinkurl{Allender.allender_polylog_genus_in_ACC0}.  The source
convention is the one explicitly defined by Hansen: fan-in-two AND/OR, unary
COPY, and literal/constant inputs.  Allender's prize question is introduced as
the unresolved polylogarithmic-genus extension of precisely Hansen's planar
characterization; Section~\ref{sec:model} records this source lineage and the
remaining wording caveat explicitly.  The Lean development is nevertheless a
\emph{partial formal verification of the mathematical theorem}: the project
treats ordinary orientable genus through an explicit definition boundary and
four named topological facts, and it imports the forward direction of Hansen's
theorem as one exact external theorem.  Those external ingredients are not
re-proved from foundational definitions in Lean.  In addition, the paper
theorem allows an arbitrary finite source circuit at $n=0$, whereas Lean's
normalized source-size predicate constrains that length as well.  The finite
wrapper extending the normalized Lean hypotheses to this broader paper-level
$n=0$ statement is not machine checked; Section~\ref{sec:lean} states the
boundary explicitly.
\end{abstract}

\tableofcontents

\section{Introduction and status}

Hansen proved that constant-width polynomial-size planar Boolean circuits
characterize $\ACC$~\cite{Hansen2006}.  Allender, Datta, and Roy proposed a
topological extension in which planarity is replaced by polylogarithmic
orientable genus~\cite{ADR2005,ECCC2004}.  The proof of their main upper-bound
theorem was later found to be incorrect; the ECCC record now explicitly notes
that the proof of Theorem 6 is wrong~\cite{ECCCComment2025}.  Allender also
listed the corresponding problem among his prize questions~\cite{Allender2023}.

The purpose of this paper is to prove the polynomial-size version of that
upper bound by a different route.  The proof never orders or pairs handles.
It uses orientable genus only through monotonicity, invariance under relabeling,
the zero-genus base case, and additivity over connected components.

\begin{theorem}[Main theorem]\label{thm:main}
Let $L\subseteq\bits^*$ be recognized by a nonuniform family
$\{C_n\}_{n\ge 0}$ of layered Boolean circuits in Hansen's standard
constant-width basis: fan-in-two AND and OR gates, unary COPY gates, input
literals and negated input literals, and Boolean constants.  Assume that there
are integers
\[
  w\ge 1,\qquad k,c\ge 0,\qquad A,B\ge 1
\]
such that, for every $n\ge 1$,
\[
  \operatorname{width}(C_n)\le w,\qquad
  |C_n|\le B(n+1)^k,\qquad
  \genus(C_n)\le A\bigl(\log_2(n+2)+1\bigr)^c.
\]
No restriction beyond finiteness is needed at the single input length $n=0$.
Then $L\in\ACC$.
\end{theorem}

The converse follows immediately from Hansen's theorem because planar graphs
have genus zero.  Hence constant-width polynomial-size circuits of
polylogarithmic orientable genus compute exactly $\ACC$.

\begin{remark}[Normalization of polynomial and polylogarithmic bounds]\label{rem:normalization}
The constants in Theorem~\ref{thm:main} are written as integers only to make
all later estimates literal.  Any usual $O((\log n)^{O(1)})$ genus bound can
be enlarged to one with an integer coefficient $A$ and integer exponent $c$.
Likewise, for $n\ge1$,
\[
 B(n+1)^k\le (n+1)^{k+B},
\]
so the coefficient $B$ may be absorbed into the polynomial exponent.  Delete
empty layers first if a presentation contains any; this does not change the
computed function or genus.  Thereafter the number of layers is at most the
number of gates.  Padding each remaining width-at-most-$w$ layer to exactly
$w$ positions therefore multiplies size by at most the constant $w$.  The new
positions can be unused Boolean-constant gates, hence isolated in the
dependency graph, so the padding does not increase genus.

The exceptional input length $n=0$ must be handled separately rather than
forced into the normalized inequality $|C_n|\le(n+1)^{k'}$.  The construction
below first produces, for all $n\ge1$, one fixed-modulus target family.  Its
circuit at length $0$ is then chosen to be the constant value $L(\epsilon)$,
using the same modulus.  This changes neither the uniform depth bound nor the
polynomial size bound and is legitimate in the nonuniform setting.

The current Lean declaration is slightly narrower at exactly this point:
\nolinkurl{CircuitFamily.PolynomialSize} requires the normalized size inequality
at every input length, including $n=0$.  Thus Lean checks the reduction for an
already normalized source family; the finite length-zero patch just described
is a paper-level wrapper and is not claimed to be machine checked.
\end{remark}

\begin{remark}[Circuit convention inherited by Open Question 3]\label{rem:allender-model}
The one-line wording of Allender's Open Question~3 does not repeat gate arities.
Its surrounding source context, however, is quite specific.  Section~4 of
Allender's prize column first recalls Hansen's characterization of $\ACC$ by
planar constant-width circuits, then describes Allender--Datta--Roy as an
extension of that result from genus zero to genus $\log^{O(1)}n$, retracts the
published proof, and immediately asks Open Question~3~\cite{Allender2023}.
Hansen explicitly defines a constant-width circuit to use fan-in-two AND/OR,
fan-in-one COPY, and literal or Boolean-constant inputs~\cite{Hansen2006}.
Allender--Datta--Roy likewise present their main theorem as an extension of
Hansen's characterization~\cite{ADR2005,ECCC2004}.  Their short Definition~2
does not restate the AND/OR arity, and later uses one-input ``dummy OR'' gates,
which are semantically COPY gates.  We therefore use Hansen's explicit
convention throughout.

This paragraph is a statement about the source convention supported by the
primary-source chain, not a topology-preserving normalization theorem.  The
paper makes no claim for a separate model in which source AND/OR gates may have
arbitrary higher fan-in.  If such a broader reading were intended, an
additional genus-preserving normalization argument would be required.
\end{remark}

\begin{remark}[Why polynomial size is part of the statement]\label{rem:size}
Allender's short prize wording refers to the question arising from the
retracted upper bound of Allender--Datta--Roy.  Their Theorem~6 and the ECCC
abstract state the claim for constant-width circuits of polynomial size and
polylogarithmic genus.  Theorem~\ref{thm:main} proves precisely that
polynomial-size formulation.  Without a polynomial-size restriction, a
sufficiently long constant-width sequential circuit can compute arbitrary
Boolean functions.  If the short prize wording were instead read literally as
omitting every size restriction, that would be a different, strictly stronger
statement and is not claimed here.
\end{remark}

\begin{verification}[Partial Lean verification and trust boundary]\label{ver:scope}
The public repository
\url{https://github.com/Grisha-Pochuev/allender-polylog-genus-lean}
contains a machine-checked implementation of the new reduction in a concrete
circuit model.  For reproducibility, the proof source referenced by this
manuscript is the exact commit
\nolinkurl{37f90d350278a40c360375c7f8731c46a2610ec5}, not a moving branch name.
At that proof state the final declaration is\par\smallskip\noindent
\nolinkurl{Allender.allender_polylog_genus_in_ACC0}.\par\smallskip
GitHub pull-request workflow run \texttt{31135088313} completed successfully
with that recorded PR-head SHA.  GitHub checked the corresponding generated
pull-request merge result, rejected \texttt{sorry}/\texttt{admit}, completed
\texttt{lake build}, compiled the axiom audit, and replayed every built project
module with \texttt{leanchecker}.  Later documentation/workflow changes on the development branch
\nolinkurl{formalization/canonical-components-v2} are not used as the
proof-version identifier here.

This should be described as a \emph{partial formal verification of the
mathematical theorem}, not as a foundational verification of every imported
result.  The development leaves an explicit external boundary: the symbol
representing ordinary orientable genus, four standard properties of that
invariant, and the forward family-level direction of Hansen's theorem.  Lean
checks the separator-based reduction relative to that boundary and for the
normalized source-family interface used in the repository.  The exceptional
length-zero wrapper in Remark~\ref{rem:normalization} is outside the checked
final declaration.  Section~\ref{sec:lean} records the scope exactly.
\end{verification}

\section{Circuit model and compatibility with Hansen}\label{sec:model}

We work with the layered model of Allender--Datta--Roy: a circuit has layers
\[
  V_0,V_1,\ldots,V_\ell,
\]
and every dependency edge goes from $V_i$ to $V_{i+1}$~\cite{ADR2005,ECCC2004}.
Thus no edge can jump over an intermediate layer; this exact adjacency
condition is essential to the separator argument below.  The width is
$\max_i |V_i|$.  The source gate basis used in the formalization is the one
needed by Hansen's planar theorem: fan-in-two AND and OR, unary COPY, input
literals, negated input literals, and Boolean constants.  Inputs may be used on
arbitrary layers and may be repeated.  The genus of a circuit is the ordinary
orientable genus of its underlying undirected dependency graph.

Every layer is represented using exactly $w$ positions, padding unused
positions by dummy zero values.  Thus the state space is the fixed finite set
\[
  Q:=\bits^w,
  \qquad |Q|=2^w.
\]
The crucial point is that $|Q|$ is independent of $n$.

\subsection{Gate-basis convention and fixed modulus}

We use exactly the constant-width circuit convention made explicit by Hansen:
AND and OR gates have fan-in two, COPY gates have fan-in one, and input nodes
are literals or Boolean constants~\cite{Hansen2006}.  This is also the source
gate basis implemented by the Lean development.  Remark~\ref{rem:allender-model}
explains why this is the convention inherited by the prize question from the
Hansen--Allender--Datta--Roy source chain.  In particular, the unary dummy OR
gates appearing in Allender--Datta--Roy play the role of COPY gates.

We do \emph{not} use a topology-preserving conversion from an arbitrary
higher-fan-in presentation.  Such a conversion would have to respect the
rotation of incident edges in a surface embedding and is not needed for the
theorem proved here.  Thus no local vertex-splitting lemma is hidden in the
reduction.

We also fix notation for the target class.  In the nonuniform setting used
here,
\[
  \ACC=\bigcup_{m\ge2}\ACm{m}.
\]
Membership in \(\ACm{m}\) means that one modulus \(m\) is fixed for the entire
circuit family, over all input lengths.  The existential choice of this single
modulus is made only once.  This quantifier is important in
Section~\ref{sec:hansen}; proving unrelated block families to lie somewhere in
\(\ACC\) would not by itself provide a common modulus.  We freely use Boolean
negation on already computed target bits.  This is the standard equivalent
closure convention for \(\ACm{m}\); it does not require changing the modulus
and can be implemented with constant additional depth.

\section{The layer planarization theorem}\label{sec:planarization}

Only standard global facts about orientable genus are needed:
\begin{enumerate}[label=(G\arabic*)]
  \item $\genus(H)\le \genus(G)$ when $H$ is a subgraph of $G$;
  \item injective relabeling and adding isolated vertices do not change genus;
  \item an edgeless graph has genus zero;
  \item genus is additive over connected components.
\end{enumerate}
The additivity theorem is classical~\cite{BHKY1962}.  In particular, if
$\genus(G)\le g$, then $G$ has at most $g$ nonplanar connected components.
Indeed, each nonplanar component has positive integer genus and their genera
sum to the genus of $G$.

\begin{lemma}[Weighted median layer]\label{lem:median}
Let $K$ be a nonempty connected vertex set in a layered graph.  There is a
layer index $m$ such that
\[
 \left|K\cap\bigcup_{i<m}V_i\right|\le \frac{|K|}{2},
 \qquad
 \left|K\cap\bigcup_{i>m}V_i\right|\le \frac{|K|}{2}.
\]
\end{lemma}

\begin{proof}
Scan the layer weights $|K\cap V_i|$ from left to right and take the first
index $m$ where the cumulative weight through layer $m$ reaches at least
$|K|/2$.  By minimality, the weight strictly to the left of $m$ is less than
$|K|/2$.  Since the cumulative weight through $m$ is at least $|K|/2$, the
weight strictly to the right is at most $|K|/2$.
\end{proof}

\begin{lemma}[Connected sets avoiding a layer lie on one side]\label{lem:one-side}
Let $H$ be a nonempty connected vertex set in a layered graph, where
``connected'' means that the subgraph induced by $H$ is connected.  If
$H\cap V_m=\varnothing$ for some layer index $m$, then either
\[
  H\subseteq \bigcup_{i<m}V_i
  \qquad\text{or}\qquad
  H\subseteq \bigcup_{i>m}V_i.
\]
\end{lemma}

\begin{proof}
Suppose instead that $H$ contains one vertex below $m$ and one vertex above
$m$.  Connectivity supplies a path between them using only vertices of $H$.
Every edge changes the layer index by exactly one, so that path must contain a
vertex of $V_m$, contradicting $H\cap V_m=\varnothing$.
\end{proof}

\begin{theorem}[Whole-layer planarization]\label{thm:layer-planarization}
Let $G$ be a finite layered graph with $N$ vertices and orientable genus at
most $g$.  There exists a set $J$ of layer indices with
\[
  |J|\le g\bigl(\ceil{\log_2 N}+1\bigr)
\]
such that deleting all vertices in layers indexed by $J$ leaves a planar
graph.  For $N=0$ the statement is immediate; below assume $N>0$.
\end{theorem}

\begin{proof}
Define the remainder $G_0:=G$.  At round $t$, let $\cK_t$ be the set of the
\emph{actual nonplanar connected components} of $G_t$.  Since $G_t$ is a
subgraph of $G$,
\[
  \genus(G_t)\le g.
\]
Genus additivity therefore gives $|\cK_t|\le g$.

For each $K\in\cK_t$, choose a weighted median layer $m_t(K)$ as in
Lemma~\ref{lem:median}.  Let
\[
  J_t:=\{m_t(K):K\in\cK_t\}
\]
and obtain $G_{t+1}$ by deleting from $G_t$ every vertex in every layer from
$J_t$.  Distinct components may choose the same global layer, so
\[
  |J_t|\le |\cK_t|\le g.
\]

We now make explicit the parent step that is easy to leave implicit.  Let $H$
be any nonplanar connected component of $G_{t+1}$.  Since $G_{t+1}$ is a
subgraph of $G_t$, the connected set $H$ is contained in a unique connected
component $K$ of $G_t$.  Because $H$ is nonplanar and $H\subseteq K$, genus
monotonicity implies that $K$ is nonplanar.  Thus $K\in\cK_t$: every
next-round nonplanar component has a canonical nonplanar parent.

The parent $K$ selected its median layer $m_t(K)$, and this whole layer was
deleted.  Hence $H$ avoids $V_{m_t(K)}$.  Although the round deletes all
layers in $J_t$ simultaneously, this causes no extra hypothesis: $H$ is still
a connected vertex set in $G_t$ and it avoids its parent's chosen median
layer.  Lemma~\ref{lem:one-side} therefore places $H$ wholly below or wholly
above that layer.  Since $H\subseteq K$, the median property gives
\[
  |H|\le \frac{|K|}{2}.
\]

Consequently, if a nonplanar component survives to round $r$, repeatedly taking
canonical parents produces a chain
\[
  K_0\supseteq K_1\supseteq\cdots\supseteq K_r
\]
with $K_t\in\cK_t$ and
\[
  |K_r|\le \frac{N}{2^r}.
\]
Take $r=\ceil{\log_2 N}+1$.  Then $N/2^r<1$, contradicting the nonemptiness of
$K_r$.  Hence $G_r$ has no nonplanar connected component.  By genus
additivity, every connected component has genus zero and therefore
$\genus(G_r)=0$, so $G_r$ is planar.

Finally let $J=\bigcup_{t<r}J_t$.  Then
\[
  |J|\le \sum_{t<r}|J_t|\le gr
  =g\bigl(\ceil{\log_2 N}+1\bigr).
\]
Deleting $J$ from $G$ produces precisely the final remainder $G_r$.
\end{proof}

\begin{corollary}[Polylogarithmic cut set]\label{cor:cuts}
If $|C_n|\le(n+1)^k$ and
$\genus(C_n)\le A(\log_2(n+2)+1)^c$, then $C_n$ can be made planar by deleting
\[
  O\bigl((\log n)^{c+1}\bigr)
\]
whole layers.
\end{corollary}

\begin{proof}
If $|C_n|=0$ there is nothing to cut.  Otherwise
Theorem~\ref{thm:layer-planarization} gives at most
\[
 A(\log_2(n+2)+1)^c
 \bigl(\ceil{\log_2|C_n|}+1\bigr)
\]
cuts.  Since $|C_n|\le(n+1)^k$, the second factor is
$O(\log(n+2))$.
\end{proof}

\section{Canonical macroblocks and exact circuit semantics}\label{sec:macroblocks}

Fix $C_n$ and a planarizing set of layer indices $J$.  A transition from layer
$i$ to layer $i+1$ is called \emph{bad} if $i\in J$ or $i+1\in J$.  Every cut
layer is incident with at most two transitions, so the number $b$ of bad
transitions satisfies
\[
  b\le 2|J|.
\]
Partition the transition list canonically into singleton bad transitions and
maximal runs of nonbad transitions.  The number $M$ of macroblocks obeys
\[
  M\le b+(b+1)\le 4|J|+1.
\]
Thus $M=O((\log n)^{c+1})$.

Every good macroblock contains no cut layer.  Its dependency graph is literally
a subgraph of the planar remainder, hence is planar.  Every bad macroblock is
a single transition between two width-$w$ layers.

\subsection{Standalone block circuits}

For exact semantics it is useful not to replace a block's incoming boundary by
informal ``fixed values.''  Instead, a standalone block circuit receives
\emph{two kinds of inputs}:
\begin{enumerate}
  \item the original $n$ input bits $x$;
  \item $w$ boundary bits $z\in Q$ giving the state entering the block.
\end{enumerate}
The block then evaluates exactly the gates in its layer interval.  For a good
macroblock, the added boundary input layer can be drawn as a local relabeling
and set of boundary stubs adjacent to the first block layer.  Equivalently,
its dependency graph embeds as a subgraph of the ambient good-block graph after
injective relabeling and adding isolates.  Hence the standalone good-block
circuit is planar.

This is the precise construction implemented by the Lean files
\texttt{MacroblockCircuit.lean} and \texttt{BlockCircuit.lean}.

\subsection{Fixed-boundary block circuits}\label{subsec:fixed-boundary}
The $(n+w)$-input circuit is useful for stating block semantics and for the
planarity comparison, but it is \emph{not} the circuit to which Hansen is
ultimately applied.  For every fixed incoming state $p\in Q$ and every output
coordinate $j\in\{1,\ldots,w\}$, define an ordinary $n$-input circuit
$P_{\cB,p,j}$ as follows: its first layer consists of $w$ Boolean constants
encoding $p$, and the remaining layers are exactly the layers of the
macroblock; coordinate $j$ of the last layer is the output.

\begin{lemma}[Fixed-boundary specialization]\label{lem:fixed-boundary}
For every macroblock $\cB$, state $p\in Q$, output coordinate $j$, and input
$x\in\bits^n$,
\[
  P_{\cB,p,j}(x)
  = \bigl(\text{state produced by $\cB$ from $p$ on $x$}\bigr)_j.
\]
If the standalone $(n+w)$-input circuit for $\cB$ is planar, then
$P_{\cB,p,j}$ is planar as well.  Its size differs by at most a constant
factor depending only on $w$.
\end{lemma}

\begin{proof}
The semantic statement follows because the constant first layer evaluates to
$p$, after which the original macroblock layers are evaluated unchanged.
For planarity, compare the fixed-boundary circuit with the standalone circuit
of the preceding subsection.  Their layer counts and all dependency-parent
sets are identical: in the first layer both a boundary-input gate and a
constant gate have no parents, and every later gate has exactly the same
parents in the preceding width-$w$ layer.  Hence their underlying dependency
graphs are identical after the obvious layer-position relabeling.  Planarity
of the standalone graph therefore implies planarity of the fixed-boundary
graph.  This is the construction checked in
\texttt{FixedBoundaryCircuit.lean}.
\end{proof}

\subsection{Relations on the fixed state space}

For a block $\cB$ and states $p,q\in Q$, define
\[
 R_{\cB}(p,q;x)=1
\]
when the standalone block circuit, run on original input $x$ and incoming
boundary state $p$, produces outgoing state $q$.

Let $I_p(x)$ assert that the actual first circuit layer has state $p$, and let
$A_q$ assert that the designated output bit is one in the final state $q$.
If $\cB_1,\ldots,\cB_M$ are the canonical macroblocks, then the original
circuit accepts exactly when
\begin{equation}\label{eq:acceptance}
  \bigvee_{p_0,p_M\in Q}
  I_{p_0}(x)\wedge
  (R_{\cB_1}\circ\cdots\circ R_{\cB_M})(p_0,p_M;x)
  \wedge A_{p_M}
\end{equation}
is true.

This equality is not merely a heuristic decomposition.  A deterministic
circuit evaluation gives a unique state after every layer; restricting that
state sequence to block boundaries gives the forward direction of
\eqref{eq:acceptance}.  Conversely, the relation of each block is the exact
evaluation relation of that block, and adjacent macroblocks partition the
entire transition list without gaps or overlaps.  Therefore compatible block
witnesses concatenate to the unique full evaluation.  The corresponding Lean
statements are
\nolinkurl{flatten_canonicalMacroblockLayers_eq_tail},
\nolinkurl{compose_macroblockRelations_eq_tailSegment}, and
\nolinkurl{accept_cons_iff_macroblockRelations}.

\section{One common Hansen simulation for all good blocks}\label{sec:hansen}

A major logical pitfall is the following.  From ``each block lies in $\ACC$''
one may not conclude that all blocks are simultaneously available in one
fixed class $\ACm{m}$, because the modulus could in principle depend on the
block family.  We avoid this by invoking Hansen only once, on one combined
planar family.

\begin{lemma}[Simultaneous family padding]\label{lem:padding}
Fix constants $d,k,w_0$.  Suppose that for each $n\ge 1$ there are at most $(n+1)^d$
planar width-$w_0$ circuits
\[
  P_{n,0},\ldots,P_{n,T(n)-1}
\]
of size at most $(n+1)^k$.  Then all these circuits admit simulations by
$\ACm{m}$ circuits with one modulus $m$, one constant depth bound, and one
polynomial size bound independent of $n$ and $t$.
\end{lemma}

\begin{proof}
Choose the code length
\[
  N(n,t):=(n+1)^{d+2}+t,
  \qquad 0\le t<T(n).
\]
The ranges for consecutive $n$ are disjoint because the gap
\[
  (n+2)^{d+2}-(n+1)^{d+2}
\]
is larger than $(n+1)^d$, while $t<(n+1)^d$.  Finitely many exceptional input lengths may be hardwired.  At coded length $N(n,t)$ place
$P_{n,t}$, viewed as a circuit with additional ignored inputs; use a constant
circuit at uncoded lengths.  This gives a single constant-width planar family
of polynomial size as a function of the new input length.

Apply Hansen's family-level theorem once to this combined family.  By the
definition of $\ACC$, the result supplies one fixed modulus $m$, one depth
bound, and one polynomial size bound for the entire family.  To recover an
$n$-input simulation of $P_{n,t}$ from the target circuit at coded length
$N(n,t)$, fix every added input coordinate to zero and replace target input
gates reading those fixed coordinates by the corresponding Boolean constants.
Because the padded source function ignores those coordinates, the restricted
target computes exactly $P_{n,t}$.  This restriction changes neither the
modulus nor the depth and does not increase size.  Since $N(n,t)$ is polynomial
in $n$, the resulting size is polynomial in $n$.
\end{proof}

For every good macroblock $\cB$, every fixed incoming state $p\in Q$, and
every output coordinate $j$, use the \emph{$n$-input fixed-boundary circuit}
$P_{\cB,p,j}$ from Lemma~\ref{lem:fixed-boundary}.  These circuits are planar,
have the original input length $n$, and have polynomial size.  There are only
\[
  M\,|Q|\,w=O((\log n)^{c+1})
\]
such circuits at length $n$, hence polynomially many.  Package exactly this
batch into the single padded family of Lemma~\ref{lem:padding}, and apply
Hansen once.  We obtain one common modulus $m$, one common depth bound, and one
common polynomial size bound for every good-block/state/output query.

For fixed $p,q\in Q$, the predicate $R_{\cB}(p,q;x)$ is the conjunction, over
the $w$ output coordinates, of either $P_{\cB,p,j}(x)$ or its complement,
according to the corresponding bit of $q$.  Negations are therefore applied
only to already obtained $\ACm{m}$ output bits.  Because $w$ is constant, this
adds only constant depth and polynomial size.  No free boundary input remains:
$p$ has been fixed by constants \emph{before} the Hansen family is formed.
This is exactly the separation between \texttt{BlockCircuit.lean},
\texttt{FixedBoundaryCircuit.lean}, and \texttt{GoodBlockBatch.lean} in the
formal development.

A bad macroblock is one transition between two constant-size states.  For
fixed $p,q$ its relation depends on at most a constant number of input literals,
so it has a constant-size $\mathrm{AC}^0$ implementation.  Pad these circuits
to the same depth.  We have proved:

\begin{proposition}[Common block-relation complexity]\label{prop:block-relations}
There is one fixed modulus $m$, one fixed depth $D$, and a constant $a$ such
that every entry $R_{\cB}(p,q;x)$ of every macroblock relation is computed by
an $\ACm{m}$ circuit of depth at most $D$ and size at most $(n+1)^a$.
\end{proposition}

\section{Constant-depth composition of all macroblocks}\label{sec:composition}

Let
\[
  \ell_n:=\log_2(n+2)+1.
\]
From the cut bound and $|C_n|\le(n+1)^k$, there is a constant $C_0$ depending
only on $A,k$ such that the number $M(n)$ of macroblocks satisfies
\begin{equation}\label{eq:block-bound}
  M(n)\le (C_0\ell_n)^{c+1}.
\end{equation}
For example, the formal proof uses the concrete coefficient
$C_0=4A(k+1)+1$ after the harmless integral-log conventions built into the
Lean definitions.

Set the integer block length
\[
  L(n):=\ceil{C_0\ell_n}.
\]
Then $L(n)\ge C_0\ell_n\ge1$, so \eqref{eq:block-bound} implies
$M(n)\le L(n)^{c+1}$.  Consider a consecutive group of the integer number
$s\le L(n)$ of relations $R_1,\ldots,R_s$.  For endpoints $p,q\in Q$, their
composite entry is exactly
\begin{equation}\label{eq:explicit-comp}
 \bigvee_{u_1,\ldots,u_{s-1}\in Q}
 \bigwedge_{i=1}^{s} R_i(u_{i-1},u_i;x),
 \qquad u_0=p,\ u_s=q.
\end{equation}
The number of candidate intermediate sequences is at most
\[
 |Q|^{L(n)-1}
 \le (2^w)^{C_0(\log_2(n+2)+1)+1}
 =(n+2)^{O(wC_0)},
\]
which is polynomial in $n$; the ceiling contributes only a constant factor.
Hence one grouping round adds one unbounded AND layer and one unbounded OR
layer and only polynomially many gates.

The size accounting can be made explicit.  If every entry circuit entering a
round has size at most $(n+1)^a$, then one group has at most
$|Q|^{L(n)}$ trajectories, each trajectory conjoins at most $L(n)$ entry
circuits.  Thus, up to a constant depending on $|Q|^2$ for all endpoint pairs,
its size is bounded by
\[
  |Q|^{L(n)}L(n)(n+1)^a+(n+1)^{O(1)}
  =(n+1)^{a+O(1)}.
\]
The exponent increase is a constant because $w,C_0$ are fixed.  After a
constant number of rounds the result is still polynomial size.

Each round replaces at most $L(n)$ consecutive relations by their composite.
Since $M(n)\le L(n)^{c+1}$ and $L(n)$ is a positive integer, after one round
the number of groups is at most $L(n)^c$; iterating this elementary ceiling
bound leaves at most one relation after $c+1$ rounds.  The Lean development
proves the corresponding exact finite-list statement, including incomplete
final groups and the zero/one-relation boundary cases.
The depth increases by only a constant per round, so the final depth is
constant because $c$ is fixed.  No new modular gate is introduced, hence the
same modulus $m$ is preserved throughout.

We therefore obtain:

\begin{lemma}[Finite-state polylogarithmic composition]\label{lem:composition}
Under Proposition~\ref{prop:block-relations} and the macroblock bound
\eqref{eq:block-bound}, every entry of
\[
  R_{\cB_1}\circ\cdots\circ R_{\cB_M}
\]
is computed by a polynomial-size constant-depth $\ACm{m}$ family for the same
fixed modulus $m$.
\end{lemma}

\section{Proof of the main theorem}\label{sec:mainproof}

\begin{proof}[Proof of Theorem~\ref{thm:main}]
Work in the Hansen-compatible source model of Section~\ref{sec:model}.  We first
reserve the exceptional length $n=0$ for the final finite patch described
below, and hence assume $n\ge1$.  By Corollary~\ref{cor:cuts}, delete
$O((\log n)^{c+1})$ whole layers so that the remaining dependency graph is
planar.  Section~\ref{sec:macroblocks} partitions the circuit into
$M=O((\log n)^{c+1})$ canonical macroblocks, each either a planar good block or
a single bad transition.

Section~\ref{sec:hansen} constructs, using one combined padded planar family,
one fixed modulus $m$ and uniform constant-depth polynomial-size $\ACm{m}$
circuits for every macroblock-relation entry.  Section~\ref{sec:composition}
composes all macroblock relations in $c+1$ logarithmic blocking rounds while
keeping the same modulus, constant depth, and polynomial size.

Finally, the exact acceptance identity \eqref{eq:acceptance} adds only a
constant-size Boolean wrapper because $|Q|=2^w$ is fixed.  Thus, for every
$n\ge1$, the construction gives circuits in one and the same $\ACm{m}$ family,
with a common constant-depth bound and a common polynomial size bound.  At
input length $0$, choose instead a constant target circuit with value
$L(\epsilon)$, using the same modulus $m$.  This finite nonuniform modification
preserves the depth and size bounds.  The resulting family recognizes exactly
$L$, and therefore $L\in\ACC$.
\end{proof}

\begin{corollary}\label{cor:equality}
Constant-width polynomial-size circuits of polylogarithmic orientable genus
compute exactly $\ACC$.
\end{corollary}

\begin{proof}
Theorem~\ref{thm:main} gives the upper bound.  Hansen's theorem gives a
constant-width polynomial-size planar circuit family for every $\ACC$ language,
and planar graphs have genus zero~\cite{Hansen2006}.
\end{proof}

\section{Why the earlier handle argument is not used}\label{sec:oldproof}

The invalid argument in Allender--Datta--Roy analyzed an embedding by cutting
and organizing handles and relied on a local ordering assertion about active
handle connections.  Allender reports a counterexample to that assertion, and
the ECCC record marks the proof of the main theorem as incorrect
\cite{Allender2023,ECCCComment2025}.

The proof above has no corresponding step.  It never chooses an embedding,
orders handles, pairs attachment arcs, or asserts that handles lie in
longitudinal stripes.  Topology enters only through global genus inequalities.
The actual cutting operation is combinatorial: delete whole circuit layers.
The halving argument is then a statement about connected components in a
layered graph.

\section{Partial Lean verification and exact trust boundary}\label{sec:lean}

The partial formal verification is maintained in the public repository
\url{https://github.com/Grisha-Pochuev/allender-polylog-genus-lean}.  This
manuscript pins the proof source to commit
\nolinkurl{37f90d350278a40c360375c7f8731c46a2610ec5}; a moving branch name is
not used as the reproducibility identifier.  Within the project's concrete
circuit model and explicit external theorem boundary, the final checked
declaration at that proof state has the form
\begin{verbatim}
theorem allender_polylog_genus_in_ACC0 (F : CircuitFamily)
    (hsourceSize : F.PolynomialSize)
    (hsourceGenus : F.PolylogGenus) : InACC0 F.language
\end{verbatim}
Fixed width is structural in \texttt{CircuitFamily}, and the source gate
constructors are fan-in-two AND/OR, unary COPY, literals, and constants.  The
predicate \texttt{F.PolynomialSize} requires its normalized size inequality at
every input length, including $n=0$.  Therefore this final declaration is
literally a theorem about the normalized repository interface; the finite
length-zero patch used to pass from Theorem~\ref{thm:main} to that interface is
not part of the checked declaration.

GitHub pull-request workflow run \texttt{31135088313}, whose recorded PR-head
SHA is\par\noindent
\nolinkurl{37f90d350278a40c360375c7f8731c46a2610ec5},\par\noindent
completed successfully.  GitHub Actions checked its generated pull-request
merge result, used Lean 4.32.2 and mathlib 4.32.2, and performed all of the
following:
\begin{enumerate}
  \item rejected \texttt{sorry} and \texttt{admit};
  \item ran \texttt{lake build} over the imported project;
  \item compiled \texttt{Allender/AxiomAudit.lean};
  \item replayed every built project module with \texttt{leanchecker}.
\end{enumerate}
All four stages completed successfully.

\subsection{What is proved internally}

The checked development includes, among other items:
\begin{itemize}
  \item concrete source-circuit semantics and the fixed state space $Q$;
  \item whole-layer deletion and the fact that surviving paths cannot cross a
        deleted layer;
  \item weighted median existence and the halving of every connected child;
  \item actual nonplanar remainder components, canonical parents, and the
        unconditional layer-planarization theorem inside the graph model;
  \item the canonical macroblock partition and its exact relation semantics;
  \item standalone good-block circuits and their planarity;
  \item the padding construction that packages all good blocks into one
        family before Hansen is used;
  \item concrete $\ACm{m}$ syntax and semantics;
  \item finite-state relation composition, exact semantics of every blocking
        round, and polynomial size/depth bounds;
  \item the final target family, exact recognition, constant depth, and
        polynomial size.
\end{itemize}

\subsection{What remains external}

The proof assistant does not rebuild orientable surface embeddings or
Hansen's paper from first principles.  Instead, its trust boundary is visible
in the axiom audit.  It consists of one external \emph{definition boundary}
for orientable genus, four standard facts about that invariant, and one exact
external theorem of Hansen.  Thus it is more precise to speak of six external
interfaces than of six independently formalized theorems.

\paragraph{Orientable genus.}
\texttt{Allender/OrientableGenus.lean} contains five external declarations:
\begin{enumerate}
  \item \texttt{genus}: an external definition boundary intended to denote
        ordinary orientable graph genus as a natural number;
  \item \nolinkurl{genus_mono}: monotonicity under spanning subgraphs;
  \item \nolinkurl{genus_map}: invariance under injective relabeling and added
        isolates;
  \item \nolinkurl{genus_bot}: an edgeless graph has genus zero;
  \item \nolinkurl{genus_eq_sum_components}: additivity over connected
        components~\cite{BHKY1962}.
\end{enumerate}
No separator or planarization theorem is assumed.

\paragraph{Hansen.}
\texttt{Allender/Hansen.lean} contains one external family-level theorem,
\nolinkurl{Hansen.planar_constantWidth_polySize_to_ACC0}, corresponding to
the forward direction of Hansen's published characterization
\cite{Hansen2006}.  It is stated for the concrete source-family model and the
concrete target predicate \texttt{InACC0}; it is not an arbitrary per-block
simulation axiom.

Accordingly, the precise machine-checked claim is: \emph{relative to this
explicit genus interface and the exact external Hansen theorem, Lean checks
the complete separator-based reduction implemented in the repository, from
the normalized fan-in-two source-family hypotheses to construction of a
fixed-modulus $\ACC$ family.}  The paper-level finite correction at $n=0$ is
not part of that declaration.  Because the external interfaces themselves are
not derived from foundational definitions in the project, the overall theorem
should be described as only partially formalized.

\section{Adversarial checks addressed in this version}\label{sec:audit}

This version incorporates the main failure modes identified during an
adversarial audit.

\begin{enumerate}[leftmargin=2.2em]
  \item \textbf{Existence of the nonplanar ancestor chain.}
  The proof of Theorem~\ref{thm:layer-planarization} explicitly assigns every
  next-round nonplanar component to the unique connected parent containing it,
  and proves that this parent is itself nonplanar.

  \item \textbf{Simultaneous deletion of several median layers.}
  Lemma~\ref{lem:one-side} is stated for every connected vertex set avoiding
  the chosen layer, not only for a component produced by deleting that one
  layer in isolation.  Therefore the halving argument remains literal when all
  median layers in $J_t$ are deleted simultaneously.

  \item \textbf{Compatibility with Hansen's gate model and the prize source.}
  Remark~\ref{rem:allender-model} records the primary-source lineage: Allender
  introduces Open Question~3 as the unresolved polylogarithmic-genus extension
  of Hansen's planar characterization, and Hansen explicitly defines fan-in-two
  AND/OR plus unary COPY.  This version does not claim an unproved
  topology-preserving normalization from arbitrary higher-fan-in gates.

  \item \textbf{Boundary values of a macroblock.}
  The standalone $(n+w)$-input circuit is used only to expose exact semantics
  and a planar graph embedding.  Before Hansen is invoked, the boundary state
  $p$ is fixed by a constant first layer, producing a genuine $n$-input planar
  circuit $P_{\cB,p,j}$.  Thus no free boundary input is silently discarded.

  \item \textbf{One common modulus and exact de-padding.}
  Hansen is applied once to a single padded family containing every good
  block/output instance, not separately to unrelated families.  Afterward the
  artificial input coordinates are explicitly fixed to constants in the
  target circuit; no appeal is made to an implicit change of input length.

  \item \textbf{Zero bits of the target state.}
  Hansen is applied to output-bit circuits; state equality is formed afterward
  by complementing the required zero-output bits and conjoining a constant
  number of conditions.

  \item \textbf{Polynomial size and integer grouping of relation composition.}
  Section~\ref{sec:composition} uses the integer group length
  $L(n)=\ceil{C_0(\log_2(n+2)+1)}$, gives an explicit polynomial trajectory
  count including the ceiling overhead, and proves that the exact integer
  group count drops to at most one after $c+1$ rounds.

  \item \textbf{Cut layers really disappear from good blocks.}
  Both transitions incident to every cut layer are bad, so a maximal good run
  contains no cut layer at all and is a genuine subgraph of the planar
  remainder.

  \item \textbf{Polynomial-size scope and the length-zero case.}
  Remark~\ref{rem:size} identifies the polynomial-size Allender--Datta--Roy
  statement proved here and does not silently claim the stronger size-free
  reading of the short prize wording.  Remark~\ref{rem:normalization} treats
  $n=0$ by a finite target-family patch, rather than incorrectly forcing it
  into the normalized inequality $|C_n|\le(n+1)^k$, and states that this wrapper
  is outside the current Lean declaration.

  \item \textbf{Types of asymptotic constants.}
  The theorem takes $w,k,c,A,B$ to be explicit integers with the displayed
  signs.  Standard real-valued or big-$O$ formulations reduce to this form by
  increasing coefficients and taking ceilings of exponents.

  \item \textbf{Partial rather than foundational Lean verification.}
  Section~\ref{sec:lean} distinguishes the internally checked reduction from
  the external definition/theorem boundary.  In particular, a green build is
  not represented as a from-first-principles formalization of orientable
  surfaces or of Hansen's theorem.

  \item \textbf{Pinned Lean proof state.}
  The manuscript cites the exact verified proof-source commit rather than a
  moving branch tip, and distinguishes the PR-head SHA from the generated
  pull-request merge result checked by GitHub Actions.

  \item \textbf{Nonuniformity.}
  Nothing in the proof claims an efficient uniform algorithm for finding the
  cut set or generating the target family.  The theorem is nonuniform, as are
  the circuit-class statements being compared.
\end{enumerate}

\section{Quantitative summary}\label{sec:quant}

After the coefficient normalization of Remark~\ref{rem:normalization}, for $n\ge1$ write
\[
  |C_n|\le(n+1)^k,
  \qquad
  \genus(C_n)\le A(\log_2(n+2)+1)^c
  \qquad (n\ge1),
\]
and set $\ell_n=\log_2(n+2)+1$.  The planarization theorem gives at most
\[
  A\ell_n^c\bigl(\log_2 |C_n|+1\bigr)
  \le A(k+1)\ell_n^{c+1}
\]
cut layers, up to the elementary integer-log convention.  The canonical
macroblock count is at most four times the number of cuts plus one.  Thus, with
\[
  C_0=4A(k+1)+1,
\]
one may bound the number of macroblocks by $(C_0\ell_n)^{c+1}$.
With the integer group length $L(n)=\ceil{C_0\ell_n}$, one has
$M(n)\le L(n)^{c+1}$, and grouping at most $L(n)$ consecutive relations per
round leaves at most one relation after $c+1$ rounds.

If the common block simulator has depth $D$, the final depth is
$D+O(c)$; the formal development tracks a concrete constant.  The size
exponent is not optimized.  It depends only on the fixed width, the source
polynomial exponent, the genus exponent and coefficient, and the polynomial
overhead in Hansen's theorem.

\section{Correspondence with the Lean development}\label{sec:map}

The following map is intended to make independent checking easier.

\small
\begin{longtable}{>{\raggedright\arraybackslash}p{0.30\textwidth} >{\raggedright\arraybackslash}p{0.62\textwidth}}
\toprule
Paper component & Main Lean files / declarations \\
\midrule
\endhead
Circuit semantics and fixed state space &
\texttt{Gate.lean}, \texttt{CircuitLayer.lean}, \texttt{Circuit.lean},
\texttt{CircuitFamily.lean}, \texttt{FiniteState.lean} \\
\addlinespace
Layer separation and weighted medians &
\texttt{LayeredGraph.lean}, \texttt{LayeredWalk.lean},
\texttt{FiniteComponent.lean}, \texttt{MedianExistence.lean} \\
\addlinespace
Canonical nonplanar parents and planarization &
\texttt{CanonicalComponents.lean}, \texttt{CanonicalPlanarization.lean},
\nolinkurl{exists_planarizing_layer_set} \\
\addlinespace
Macroblock partition, exact semantics, and fixed boundaries &
\texttt{MacroblockPartition.lean}, \texttt{MacroblockCircuit.lean},
\texttt{BlockCircuit.lean}, \texttt{FixedBoundaryCircuit.lean},
\nolinkurl{accept_cons_iff_macroblockRelations} \\
\addlinespace
Hansen interface and simultaneous padding &
\texttt{Hansen.lean}, \texttt{Padding.lean}, \texttt{InputPadding.lean},
\texttt{SimultaneousHansen.lean}, \texttt{GoodBlockBatch.lean} \\
\addlinespace
Relation circuits and repeated composition &
\texttt{FiniteRelationComposition.lean},
\texttt{MacroblockCompositionRounds.lean},
\texttt{RelationCompositionRounds.lean}, \texttt{PolynomialBounds.lean} \\
\addlinespace
End-to-end target family &
\texttt{AcceptanceCircuit.lean}, \texttt{MainTheorem.lean},
\nolinkurl{allender_polylog_genus_in_ACC0} \\
\addlinespace
Trusted-dependency report & \texttt{AxiomAudit.lean} \\
\bottomrule
\end{longtable}
\normalsize

\section{Conclusion}

The obstruction in the earlier proof was tied to the geometry of handle
attachments.  The separator argument bypasses that geometry.  Low genus
bounds the number of nonplanar connected components at every stage; layeredness
provides a whole-layer separator that halves every nonplanar descendant; and
constant width turns the resulting pieces into a finite-state composition
problem.  Hansen's planar theorem then supplies the only circuit-complexity
input needed for the good blocks.

The result is Theorem~\ref{thm:main}, stated in Hansen's explicit fan-in-two
constant-width circuit model.  The primary-source chain surrounding Allender's
Open Question~3 identifies this as the planar model whose extension to
polylogarithmic genus is at issue; the manuscript nevertheless states the
convention explicitly rather than hiding it in terminology.  The accompanying
Lean development provides a substantial partial formal verification: it checks
the new separator-based reduction in its normalized source-family interface
through construction of the final fixed-modulus family, while leaving a small,
explicit boundary for orientable genus and Hansen's published theorem.  The
finite correction at input length zero remains a paper-level wrapper.

\section*{Acknowledgments and provenance}

The project and manuscript were developed with substantial assistance from
OpenAI language models under the author's direction.  The public repository
contains the source manuscript, the Lean formalization, the axiom audit, and
reproducibility instructions.  Mathematical responsibility for submission and
for the mathematical argument, the interpretation of external interfaces, and
any claims beyond the explicitly machine-checked reduction remains with the
author.

\begin{thebibliography}{99}

\bibitem{ADR2005}
E. Allender, S. Datta, and S. Roy,
\newblock Topology Inside $\mathrm{NC}^1$,
\newblock in \emph{Proceedings of the 20th Annual IEEE Conference on Computational Complexity (CCC 2005)}, pp. 298--307, 2005.

\bibitem{ECCC2004}
E. Allender, S. Datta, and S. Roy,
\newblock Topology inside $\mathrm{NC}^1$,
\newblock ECCC Report TR04-108, 2004.
\newblock \url{https://eccc.weizmann.ac.il/report/2004/108/}

\bibitem{Allender2023}
E. Allender,
\newblock Guest Column: Parting Thoughts and Parting Shots (Read On for Details on How to Win Valuable Prizes!),
\newblock \emph{ACM SIGACT News} 54(1):63--81, March 2023.
\newblock \url{https://doi.org/10.1145/3586165.3586175}
\newblock Author draft: \url{https://people.cs.rutgers.edu/~allender/papers/sigact.news.draft.pdf}

\bibitem{ECCCComment2025}
E. Allender,
\newblock Comment to ECCC TR04-108: ``The proof of Theorem 6 is incorrect,''
\newblock accepted 2 November 2025.
\newblock \url{https://eccc.weizmann.ac.il/eccc-reports/2004/TR04-108/index.html}

\bibitem{BHKY1962}
J. Battle, F. Harary, Y. Kodama, and J. W. T. Youngs,
\newblock Additivity of the genus of a graph,
\newblock \emph{Bulletin of the American Mathematical Society} 68:565--568, 1962.

\bibitem{Hansen2006}
K. A. Hansen,
\newblock Constant Width Planar Computation Characterizes $\mathrm{ACC}^0$,
\newblock \emph{Theory of Computing Systems} 39(1):79--92, 2006.
\newblock Preliminary version: ECCC TR03-025,
\url{https://eccc.weizmann.ac.il/report/2003/025/}.

\end{thebibliography}

\end{document}
